\documentclass[11pt]{article}
\usepackage[margin=1in]{geometry}
\usepackage{amsmath,amssymb,amsthm,mathtools}
\usepackage{microtype}
\usepackage[hidelinks]{hyperref}
\setlength{\parindent}{0pt}
\setlength{\parskip}{0.3em}
\emergencystretch=2em
\allowdisplaybreaks

\newtheorem{definition}{Definition}[section]
\newtheorem{lemma}[definition]{Lemma}
\newtheorem{theorem}[definition]{Theorem}

\newcommand{\R}{\mathbb R}
\newcommand{\one}{\mathbf 1}
\newcommand{\zero}{\mathbf 0}
\newcommand{\Id}{\mathbf I}
\newcommand{\Lmat}{\mathbf L}
\newcommand{\Nmat}{\mathbf N}
\newcommand{\gam}{\boldsymbol\gamma}
\newcommand{\Ord}{\mathcal O}
\newcommand{\cof}{\operatorname{cof}}
\newcommand{\adj}{\operatorname{adj}}
\newcommand{\rank}{\operatorname{rank}}

\begin{document}
\title{Blueprint for the Lean Proof of the Order--Mass Identity,
the Aldous--Broder Theorem, and the Stopped-Forest Law}
\author{Hugh Zheng}
\date{}
\maketitle

This document states the definitions, lemmas and proofs that the Lean
development in \path{EAB/Paper/} follows, at the index types used in
Lean. The docstrings of the Lean files cite its labels, such as
\texttt{lem:orders} or \texttt{thm:stopped}.

\section{Matrix and harmonic-extension notation}

Throughout, $V$ is a finite set with $n=|V|\geq2$ and
$\Lmat\in\R^{V\times V}$. For subsets $R,C\subseteq V$,
$\Lmat[R,C]$ is the submatrix with rows $R$ and columns $C$, and
$\Lmat[U]:=\Lmat[U,U]$. All matrix and vector indices are the indicated
finite sets. We use the finite Leibniz formula for determinants
also when the index set is empty. There is exactly one permutation
of $\varnothing$, with sign $1$, and its term is an empty product,
equal to $1$. Thus the determinant of the unique $0\times0$
matrix is $1$. This value is needed only to read
Theorem~\ref{thm:order-mass} when $|B|=1$; the proof of
Theorem~\ref{thm:single-root} and Sections 5--7 do not use it.
The vectors $\one_S,\zero_S\in\R^S$ are constant, and $\mathbf e_u$
is the coordinate vector at $u$. For a nonempty finite set $S$ and a
matrix $A\in\R^{S\times S}$, write
\[
\cof_u(A):=\det\bigl(A[S\setminus\{u\}]\bigr),\qquad
\mathbf d(A):=(\cof_u(A))_{u\in S}.
\]
For $D\subseteq S$, the notation $\mathbf1_D\in\R^S$ means its
indicator. Empty sums and products are $0$ and $1$.

\begin{lemma}[Identity plus rank one]\label{lem:identity-rank-one}
For any finite set $S$ and $a,y\in\R^S$,
\[
\det(\Id_S+ay^{\mathsf T})=1+y^{\mathsf T}a.
\]
\end{lemma}
\begin{proof}
Expand the determinant by multilinearity in its columns. Terms
using $a$ in at least two columns vanish. The term using no such
column is $1$, while replacing column $j$ of $\Id_S$ by $a$
has determinant $a_j$; hence the remaining terms sum to
$\sum_{j\in S}y_ja_j$.
\end{proof}

\begin{lemma}[Cofactors of a zero-row-sum matrix]\label{lem:cofactor}
Let $S$ be nonempty and $A\in\R^{S\times S}$ satisfy
$A\one_S=\zero_S$. Put $\mathbf d=\mathbf d(A)$. Then
\[
\adj(A)=\one_S\mathbf d^{\mathsf T},\qquad
\mathbf d^{\mathsf T}A=\zero_S^{\mathsf T}.
\]
For all $a,y\in\R^S$,
\begin{equation}\label{eq:cofactor-rank-one}
\det(A+ay^{\mathsf T})
= (y^{\mathsf T}\one_S)(a^{\mathsf T}\mathbf d).
\end{equation}
\end{lemma}
\begin{proof}
The nonzero vector $\one_S$ lies in $\ker A$, so $\det A=0$.
If $\rank A=|S|-1$, then $\ker A=\R\one_S$; the identity
$A\adj(A)=0$ makes each column of $\adj(A)$ constant, with value
its diagonal entry $\cof_u(A)$ in column $u$. If
$\rank A\leq |S|-2$, every $(|S|-1)$-minor vanishes and both
$\adj(A)$ and $\mathbf d$ are zero. Thus in either case
$\adj(A)=\one_S\mathbf d^{\mathsf T}$. The identity
$\adj(A)A=\det(A)\Id_S=0$ gives
$\mathbf d^{\mathsf T}A=0$ by reading any row.

Multilinearity in columns, with terms that replace two columns by
multiples of $a$ vanishing, gives
\[
\det(A+ay^{\mathsf T})
=\det A+\sum_{j\in S}y_j\det(A\text{ with column }j
\text{ replaced by }a).
\]
Expansion of a replaced column gives
$\det(A\text{ with column }j\text{ replaced by }a)
=\sum_{u\in S}a_u\adj(A)_{ju}=a^{\mathsf T}\mathbf d$.
Since $\det A=0$, the displayed sum is
$(\sum_jy_j)(a^{\mathsf T}\mathbf d)$.
\end{proof}

\begin{definition}[Harmonic extension]\label{def:harmonic}
Let $U\subsetneq V$ and $\det\Lmat[U]\neq0$. For a boundary vector
$m\in\R^{V\setminus U}$, define $\gam^U(m)\in\R^V$ by
\begin{equation}\label{eq:harmonic}
\gam^U(m)|_{V\setminus U}=m,
\qquad
\gam^U(m)|_U
=-\Lmat[U]^{-1}\Lmat[U,V\setminus U]m.
\end{equation}
When $U=\varnothing$, the second equation is empty and
$\gam^\varnothing(m)=m$. For $w\notin U$, abbreviate
$\gam^{U,w}:=\gam^U(\mathbf e_w)$.
\end{definition}

\begin{lemma}[Characterization, linearity, and unit boundary]\label{lem:harmonic}
Under Definition~\ref{def:harmonic}, $\gam^U(m)$ is the unique
$g\in\R^V$ with $g|_{V\setminus U}=m$ and
$(\Lmat g)|_U=\zero_U$. The map $m\mapsto\gam^U(m)$ is linear.
If in addition $\Lmat\one_V=\zero_V$, then
\begin{equation}\label{eq:unity}
\gam^U(\one_{V\setminus U})=\one_V,
\qquad
\sum_{w\in V\setminus U}\gam^{U,w}=\one_V.
\end{equation}
\end{lemma}
\begin{proof}
The restriction of $\Lmat\gam^U(m)$ to $U$ is
$\Lmat[U]\gam^U(m)|_U+\Lmat[U,V\setminus U]m=0$.
If two vectors have the required boundary and row values, their
difference vanishes off $U$ and its restriction to $U$ is in the
kernel of the invertible matrix $\Lmat[U]$; hence they coincide.
The formula in \eqref{eq:harmonic} is linear in $m$. Under zero row
sums, $\one_V$ has boundary value $\one_{V\setminus U}$ and is
annihilated on $U$, so uniqueness yields the first identity in
\eqref{eq:unity}. The second follows by linearity from
$\one_{V\setminus U}=\sum_{w\notin U}\mathbf e_w$.
\end{proof}

\begin{lemma}[One-vertex harmonic update]\label{lem:one-step}
Let $U\subsetneq V$, $w\notin U$, and $U'=U\cup\{w\}\subsetneq V$.
Assume $\det\Lmat[U]\neq0$ and $\det\Lmat[U']\neq0$.
For $m\in\R^{V\setminus U'}$, let $\widetilde m\in\R^{V\setminus U}$
extend $m$ by $\widetilde m_w=0$. Then
\begin{equation}\label{eq:one-step}
\gam^{U'}(m)
=\gam^U(\widetilde m)
 +\bigl(\gam^{U'}(m)\bigr)_w\gam^{U,w}.
\end{equation}
\end{lemma}
\begin{proof}
Both sides are annihilated by the rows $U$ of $\Lmat$. They agree
off $U$: at $w$ the right side is
$0+(\gam^{U'}(m))_w\cdot1$; at every other vertex outside $U$ its
second term is zero and its first term equals $m$. Uniqueness in
Lemma~\ref{lem:harmonic} proves the equality.
\end{proof}

\section{Aggregated matrices and one-vertex transfer}

\begin{definition}[Aggregation matrix]\label{def:aggregation}
For nonempty $U\subsetneq V$ with $\det\Lmat[U]\neq0$ and
$M\in\R^{(V\setminus U)\times U}$, put
\begin{equation}\label{eq:aggregation}
\Nmat(U,M)
:=\Id_U+\Lmat[U]^{-1}\Lmat[U,V\setminus U]M
\in\R^{U\times U}.
\end{equation}
The matrix $M$ has \emph{unit row sums} if
$M\one_U=\one_{V\setminus U}$.
\end{definition}

\begin{lemma}[Harmonic columns and zero row sums]\label{lem:aggregation}
For $j\in U$, let $m_j=M_{\cdot j}\in\R^{V\setminus U}$. Then
\begin{equation}\label{eq:aggregation-column}
\Nmat(U,M)_{\cdot j}
=\mathbf e_j-\gam^U(m_j)|_U.
\end{equation}
If $\Lmat\one_V=0$ and $M$ has unit row sums, then
$\Nmat(U,M)\one_U=0$.
\end{lemma}
\begin{proof}
Equation~\eqref{eq:aggregation-column} is
\eqref{eq:harmonic} applied to the column $m_j$. Summing it over
$j\in U$ and using linearity and
$\sum_jm_j=M\one_U=\one_{V\setminus U}$ gives
$\Nmat(U,M)\one_U=\one_U-\gam^U(\one_{V\setminus U})|_U=0$
by \eqref{eq:unity}.
\end{proof}

\begin{definition}[Fold]\label{def:fold}
Let $U\subsetneq V$ be nonempty, $w\in V\setminus U$,
$U'=U\cup\{w\}\subsetneq V$, and
$M'\in\R^{(V\setminus U')\times U'}$.
For $\theta\in\R^U$ with $\sum_{j\in U}\theta_j=1$, define its
fold $M=\operatorname{fold}_{w,\theta}(M')
\in\R^{(V\setminus U)\times U}$ by
\begin{equation}\label{eq:fold}
M_{wj}=\theta_j,
\qquad
M_{vj}=M'_{vj}+M'_{vw}\theta_j
\quad(v\in V\setminus U',\ j\in U).
\end{equation}
Write $s=M'_{\cdot w}\in\R^{V\setminus U'}$ and let
$\widetilde s\in\R^{V\setminus U}$ extend $s$ by zero at $w$.
For $j\in U'$, let $\widetilde m'_j\in\R^{V\setminus U}$ similarly
extend $M'_{\cdot j}$ by zero at $w$.
\end{definition}

\begin{lemma}[Fold identities]\label{lem:fold}
If $M'$ has unit row sums, so does $M$. For $j\in U$,
\begin{equation}\label{eq:fold-column}
\widetilde m'_j=M_{\cdot j}
 -\theta_j(\widetilde s+\mathbf e_w)
\quad\text{in }\R^{V\setminus U}.
\end{equation}
\end{lemma}
\begin{proof}
The row of $M$ at $w$ sums to $\sum_j\theta_j=1$. For
$v\in V\setminus U'$, its sum is
$\sum_{j\in U}M'_{vj}+M'_{vw}\sum_{j\in U}\theta_j
=\sum_{j\in U'}M'_{vj}=1$.
At $w$, the right side of \eqref{eq:fold-column} is
$\theta_j-\theta_j=0$; at $v\neq w$ it is
$M'_{vj}+M'_{vw}\theta_j-\theta_jM'_{vw}=M'_{vj}$.
\end{proof}

\begin{lemma}[One-vertex transfer]\label{lem:transfer}
Assume $\Lmat\one_V=0$. Let $U\subseteq V$ be nonempty with
$|U|\leq n-2$, choose $w\in V\setminus U$, and put
$U'=U\cup\{w\}$. Assume
$\det\Lmat[U]\neq0$ and $\det\Lmat[U']\neq0$.
Let $M'\in\R^{(V\setminus U')\times U'}$ have unit row sums,
let $\theta\in\R^U$ satisfy $\sum_{j\in U}\theta_j=1$,
and set $M=\operatorname{fold}_{w,\theta}(M')$.
With $s$ and $\widetilde s$ as in Definition~\ref{def:fold},
\begin{equation}\label{eq:transfer}
\sum_{x\in U'}\gam^{U,w}_x\,
  \cof_x\bigl(\Nmat(U',M')\bigr)
=\sum_{x\in U}\gam^U(\widetilde s+\mathbf e_w)_x\,
  \cof_x\bigl(\Nmat(U,M)\bigr).
\end{equation}
In particular, the right side is independent of $\theta$.
\end{lemma}
\begin{proof}
For $j\in U'$, let $m'_j=M'_{\cdot j}$ and define
\[
\widehat G_{xj}:=\gam^U(\widetilde m'_j)_x
\quad(x,j\in U'),\qquad
\rho_j:=\gam^{U'}(m'_j)_w,
\qquad a:=\gam^{U,w}|_{U'}.
\]
Applying \eqref{eq:one-step} to every column, followed by
\eqref{eq:aggregation-column}, gives
\begin{equation}\label{eq:transfer-rank-one}
\Nmat(U',M')=(\Id_{U'}-\widehat G)-a\rho^{\mathsf T}.
\end{equation}
Since the rows of $M'$ sum to one, linearity of harmonic extension
and \eqref{eq:unity}, read at $w$, give
$\sum_{j\in U'}\rho_j=1$. Also
$\widehat G_{wj}=\widetilde m'_j(w)=0$ for every $j$.

Put $D:=\det(\Id_{U'}-\widehat G)$. Equation
\eqref{eq:transfer-rank-one}, Lemma~\ref{lem:aggregation}, and
\eqref{eq:cofactor-rank-one}, applied to
$A=\Nmat(U',M')$, yield
\begin{equation}\label{eq:transfer-left}
D=\sum_{x\in U'}a_x\cof_x\bigl(\Nmat(U',M')\bigr).
\end{equation}
The row $w$ of $\Id_{U'}-\widehat G$ is the coordinate row
$\mathbf e_w^{\mathsf T}$. Expanding its determinant along this row,
\begin{equation}\label{eq:transfer-drop}
D=\det\bigl((\Id_{U'}-\widehat G)[U]\bigr).
\end{equation}
For $j\in U$, use \eqref{eq:fold-column}, linearity of harmonic
extension, and \eqref{eq:aggregation-column}. They give
\begin{equation}\label{eq:transfer-block}
(\Id_{U'}-\widehat G)[U]
=\Nmat(U,M)+a'\theta^{\mathsf T},
\qquad
a':=\gam^U(\widetilde s+\mathbf e_w)|_U.
\end{equation}
The fold has unit row sums by Lemma~\ref{lem:fold}, so
Lemma~\ref{lem:aggregation} gives $\Nmat(U,M)\one_U=0$.
Applying \eqref{eq:cofactor-rank-one} again and using
$\sum_j\theta_j=1$ yields
\begin{equation}\label{eq:transfer-right}
D=\sum_{x\in U}a'_x\cof_x\bigl(\Nmat(U,M)\bigr).
\end{equation}
Equations \eqref{eq:transfer-left} and \eqref{eq:transfer-right}
are \eqref{eq:transfer}.
\end{proof}

\begin{lemma}[Transfer onto a one-vertex base]\label{lem:single-base}
Assume the hypotheses of Lemma~\ref{lem:transfer} on $\Lmat$, $U$,
$w$, $U'$ and $M'$, with $U=\{r\}$. Then
\begin{equation}\label{eq:single-base}
\sum_{x\in U'}\gam^{U,w}_x\,
  \cof_x\bigl(\Nmat(U',M')\bigr)
=\gam^U(\widetilde s+\mathbf e_w)_r .
\end{equation}
Both cofactors on the left are determinants of $1\times1$ matrices.
\end{lemma}
\begin{proof}
Use the notation of the proof of Lemma~\ref{lem:transfer}.
Equation~\eqref{eq:transfer-left} identifies the left side with
$D$. Its derivation applies Lemma~\ref{lem:cofactor} to the
$2\times2$ matrix $\Nmat(U',M')$ only, and does not involve
$\theta$. By \eqref{eq:transfer-drop}, $D$ is the determinant of
the $1\times1$ matrix $(\Id_{U'}-\widehat G)[U]$. The Leibniz
formula on a one-element index set has the single term of the
identity permutation, so $D$ is the entry of that matrix at $(r,r)$.

Let $\theta\in\R^U$ have its only coordinate equal to $1$ and let
$M=\operatorname{fold}_{w,\theta}(M')$. By Lemmas~\ref{lem:fold}
and~\ref{lem:aggregation}, $\Nmat(U,M)\one_U=0$. The index set $U$
has one element, so this says $\Nmat(U,M)_{rr}=0$. The entry at
$(r,r)$ of \eqref{eq:transfer-block} is therefore
$D=0+a'_r\theta_r=\gam^U(\widetilde s+\mathbf e_w)_r$.
No determinant on an empty index set occurs.
\end{proof}

\section{Parent maps, root indicators, and growth orders}

Let $B\subsetneq V$ be nonempty, put $Z=V\setminus B$, and let
$k=|Z|$. A map $f:Z\to V$ is a \emph{$B$-parent map} if, starting at
each $v\in Z$ and repeatedly applying $f$ while the current vertex
lies in $Z$, the sequence reaches $B$. Let $q_B(v)\geq1$ be its
first hitting time of $B$, let $f^{(t)}(v)$ denote the iterates up to
that time, and put $r_B(v)=f^{(q_B(v))}(v)$.
For a typed implementation, extend $f$ to $V$ by fixing every vertex
of $B$, and use this total map to define the iterates and the least
positive hitting time. All iterates before that time lie in $Z$;
consequently $q_B(v)$, $r_B(v)$, descendant blocks, orders, and
weights depend only on the original map $f:Z\to V$. Prove this
independence before using a total-map representation in a theorem.

\begin{definition}[Root and descendant blocks]\label{def:blocks}
For $j\in B$ and $u\in Z$, set
\[
Z_j:=\{v\in Z:r_B(v)=j\},
\qquad
D_u:=\{v\in Z:\text{for some }0\leq t<q_B(v),\ f^{(t)}(v)=u\}.
\]
The root indicator $M_f\in\{0,1\}^{Z\times B}$ is
$(M_f)_{vj}=\mathbf1_{Z_j}(v)$.
A vertex $w\in Z$ with $f(w)=i\in B$ is a \emph{child} of $i$.
\end{definition}

\begin{lemma}[Root and child partitions]\label{lem:blocks}
The sets $(Z_j)_{j\in B}$ partition $Z$, so
$M_f\one_B=\one_Z$. For each $i\in B$, the sets $D_w$ over the
children $w$ of $i$ are pairwise disjoint and partition $Z_i$.
\end{lemma}
\begin{proof}
The first hit $r_B(v)$ exists and is unique for every $v\in Z$,
giving the first partition. If $v\in Z_i$, the last non-root on
its path before the first hit $i$ is the unique
$w=f^{(q_B(v)-1)}(v)$ with $f(w)=i$; thus $v\in D_w$.
Conversely, if $v\in D_w$ and $f(w)=i\in B$, the path from $v$
reaches $w$ before any root and then reaches $i$, so $v\in Z_i$.
The last non-root before the first root is unique, proving
disjointness.
\end{proof}

\begin{definition}[Growth orders and weights]\label{def:orders}
A \emph{growth order} of $f$ is an ordering
$\tau=(\tau_1,\ldots,\tau_k)$ of $Z$ such that
\[
U_l^{\tau}:=B\cup\{\tau_1,\ldots,\tau_l\}
\quad(0\leq l\leq k),\qquad
f(\tau_l)\in U_{l-1}^{\tau}
\quad(1\leq l\leq k).
\]
Write $\Ord_B(f)$ for the set of these orders. Whenever the
indicated harmonic extensions exist, set
\[
\Pi_{B,f}(\tau):=
\prod_{l=1}^{k-1}
  \gam^{U_{l-1}^{\tau},\tau_l}_{f(\tau_{l+1})},
\qquad
\mathcal G_i(B,f):=
\sum_{\substack{\tau\in\Ord_B(f)\\ f(\tau_1)=i}}
  \Pi_{B,f}(\tau)\quad(i\in B).
\]
For $k=1$, the product is $1$.
\end{definition}

\begin{lemma}[Existence and peeling of growth orders]\label{lem:orders}
Every $B$-parent map has a growth order. Suppose $k\geq2$.
If $w\in Z$ is a child
of $i\in B$, put $B'=B\cup\{w\}$ and let $f'$ be the restriction of
$f$ to $V\setminus B'$. Then $f'$ is a $B'$-parent map and
\[
\tau'\longmapsto(w,\tau')
\]
is a bijection from $\Ord_{B'}(f')$ onto the growth orders of $f$
whose first vertex is $w$. If $k\geq2$ and
$\tau'=(\tau'_1,\ldots,\tau'_{k-1})$, then
\begin{equation}\label{eq:weight-peel}
\Pi_{B,f}(w,\tau')
=\gam^{B,w}_{f(\tau'_1)}\Pi_{B',f'}(\tau').
\end{equation}
Every proper prefix base of $\tau'$ is a proper prefix base of
$(w,\tau')$, with its prefix length increased by one.
\end{lemma}
\begin{proof}
If $k\geq2$ and $w$ is any child of a root, $B'=B\cup\{w\}$
remains proper. Every remaining parent path reaches $B'$, stopping
at $w$ if it meets $w$ before $B$; hence $f'$ is a
$B'$-parent map.

If $k=1$, its only vertex $w$ must satisfy $f(w)\in B$, since
otherwise its path would never hit $B$; hence $(w)$ is a growth
order. For $k\geq2$, follow the path of any $v\in Z$ to its first
root. Its predecessor $w\in Z$ is a child of that root.
Induction on $k$, prepending $w$ to an order for $f'$,
proves existence.

For a fixed child $w$, its parent lies in $B$. After placing $w$
first, the condition on each subsequent vertex is exactly the
growth-order condition for $f'$ and base $B'$, proving the bijection.
The first factor of the product for $(w,\tau')$ is
$\gam^{B,w}_{f(\tau'_1)}$; shifting all remaining indices by one
gives the factors of $\Pi_{B',f'}(\tau')$. The equality of the
prefix bases is
$B'\cup\{\tau'_1,\ldots,\tau'_l\}
=B\cup\{w,\tau'_1,\ldots,\tau'_l\}$.
\end{proof}

\begin{lemma}[Peeling folds the root indicator]\label{lem:forest-fold}
Under the $k\geq2$ hypotheses of Lemma~\ref{lem:orders},
let $M'=M_{f'}$ and regard
$\mathbf e_i\in\R^B$ as a vector with coordinate sum one. Then
\[
M_f=\operatorname{fold}_{w,\mathbf e_i}(M'),
\qquad
\widetilde s+\mathbf e_w=\mathbf1_{D_w}
\quad\text{in }\R^Z,
\]
where $s=M'_{\cdot w}$ is the column of the new root $w$.
\end{lemma}
\begin{proof}
The row of $M_f$ at $w$ is $\mathbf e_i^{\mathsf T}$ because
$f(w)=i$. For $v\neq w$, if its path first meets $w$ before any
vertex of $B$, its new root is $w$ and its old root is $i$;
otherwise its new and old roots coincide in $B$.
Consequently, for each $j\in B$,
$(M_f)_{vj}=(M')_{vj}+(M')_{vw}\mathbf1_{\{i\}}(j)$,
which is the fold formula. The column $s$ indicates the remaining
vertices whose paths first meet $w$; adjoining $w$ itself gives
precisely $D_w$.
\end{proof}

\section{Order--mass identity}

\begin{theorem}[Order--mass identity]\label{thm:order-mass}
Let $B\subsetneq V$ be nonempty, let $f:V\setminus B\to V$ be a
$B$-parent map, and let $\Lmat\in\R^{V\times V}$ satisfy
$\Lmat\one_V=0$ and
\begin{equation}\label{eq:minor-hypothesis}
\det\Lmat[U_l^{\tau}]\neq0
\quad\text{for every }\tau\in\Ord_B(f)
\text{ and }0\leq l\leq k-1.
\end{equation}
Then, for every $i\in B$,
\begin{equation}\label{eq:order-mass}
\mathcal G_i(B,f)
=\cof_i\bigl(\Nmat(B,M_f)\bigr)
=\det\!\left[
  \bigl(\Id_B+\Lmat[B]^{-1}\Lmat[B,Z]M_f\bigr)
    [B\setminus\{i\}]
  \right].
\end{equation}
\end{theorem}
\begin{proof}
If $|B|=1$, the right side of \eqref{eq:order-mass} is
$\det(\varnothing)=1$ by the Leibniz formula, and the identity is
Theorem~\ref{thm:single-root}. The proof of that theorem uses the
present one for a two-element root set only. Assume $|B|\geq2$.
First establish the assertion for $|B|\geq2$ as a separate
intermediate theorem. In that theorem, induct on $k=|Z|$,
simultaneously for all proper root sets
with at least two elements and their parent maps. The induction
step enlarges the root set, so it stays within this class, and
every determinant below is taken on a nonempty index set.
Lemma~\ref{lem:orders} supplies a growth
order, so \eqref{eq:minor-hypothesis} at $l=0$ makes
$\Nmat(B,M_f)$ well defined.

If $k=1$, write $Z=\{w\}$ and $i_0=f(w)\in B$. The only growth
order is $(w)$; hence $\mathcal G_i(B,f)=\mathbf1_{\{i_0\}}(i)$.
Zero row sums give
$\Lmat[B]\one_B+\Lmat[B,\{w\}]\one_{\{w\}}=0$, so
$\Nmat(B,M_f)=\Id_B-\one_B\mathbf e_{i_0}^{\mathsf T}$.
Its principal cofactor is $1$ when $i=i_0$, because deleting that
row and column leaves an identity matrix. When $i\neq i_0$, the
remaining matrix is
$\Id_{B\setminus\{i\}}-\one_{B\setminus\{i\}}
\mathbf e_{i_0}^{\mathsf T}$ and has determinant $0$ by
Lemma~\ref{lem:identity-rank-one}. This proves
\eqref{eq:order-mass}.

Suppose $k\geq2$, fix $i\in B$, and abbreviate
$N=\Nmat(B,M_f)$ and $d_x=\cof_x(N)$. Lemmas
\ref{lem:blocks} and \ref{lem:aggregation} imply $N\one_B=0$;
therefore Lemma~\ref{lem:cofactor} gives
\begin{equation}\label{eq:left-kernel}
\sum_{x\in B}d_xN_{xi}=0.
\end{equation}
Group the growth orders counted by $\mathcal G_i(B,f)$ according
to their first vertex $w$, necessarily a child of $i$. For each
such $w$, use the notation $B',f',M'$ of Lemmas~\ref{lem:orders}
and \ref{lem:forest-fold}. Lemma~\ref{lem:orders} also carries the
minor condition to $f'$: every proper prefix base of a new order
is the corresponding proper prefix base of an old order beginning
with $w$. In particular, $\det\Lmat[B']\neq0$ because an order
beginning with $w$ exists. The inductive hypothesis at each
$x\in B'$, combined with \eqref{eq:weight-peel}, gives
\begin{equation}\label{eq:induction-sum}
\mathcal G_i(B,f)
=\sum_{\substack{w\in Z\\ f(w)=i}}
   \sum_{x\in B'}\gam^{B,w}_x\,
      \cof_x\bigl(\Nmat(B',M')\bigr).
\end{equation}

For each $w$, Lemma~\ref{lem:forest-fold} identifies $M_f$ as
the fold of $M'$ along $\mathbf e_i$ and identifies
$\widetilde s+\mathbf e_w=\mathbf1_{D_w}$.
Lemma~\ref{lem:transfer} changes the inner sum in
\eqref{eq:induction-sum} into
$\sum_{x\in B}\gam^B(\mathbf1_{D_w})_x d_x$.
By Lemma~\ref{lem:blocks}, the $D_w$ for the children of $i$
partition $Z_i$. Linearity in Lemma~\ref{lem:harmonic} therefore
gives
\[
\mathcal G_i(B,f)
=\sum_{x\in B}\gam^B(\mathbf1_{Z_i})_x d_x.
\]
The $i$th column of $M_f$ is $\mathbf1_{Z_i}$, so
\eqref{eq:aggregation-column} says
$\gam^B(\mathbf1_{Z_i})|_B=\mathbf e_i-N_{\cdot i}$.
Using \eqref{eq:left-kernel},
\[
\mathcal G_i(B,f)
=\sum_{x\in B}(\mathbf e_i-N_{\cdot i})_x d_x
=d_i-\sum_{x\in B}d_xN_{xi}=d_i.
\]
The definition of $N$ gives the final displayed determinant in
\eqref{eq:order-mass}.
\end{proof}

\begin{theorem}[Single-root normalization]\label{thm:single-root}
Under the hypotheses of Theorem~\ref{thm:order-mass}, if
$B=\{r\}$, then
\[
\sum_{\tau\in\Ord_{\{r\}}(f)}\Pi_{\{r\},f}(\tau)=1.
\]
\end{theorem}
\begin{proof}
Put $B=\{r\}$, $Z=V\setminus\{r\}$ and $k=|Z|=n-1$. The proof
uses no determinant on an empty index set.

If $k=1$, write $Z=\{w\}$. By Lemma~\ref{lem:orders} the only
growth order is $(w)$, and $\Pi_{B,f}((w))=1$ by
Definition~\ref{def:orders}. The sum is $1$.

Suppose $k\geq2$. Every growth order $\tau$ has
$f(\tau_1)\in U_0^\tau=\{r\}$, so its first vertex is a child of
$r$. Lemma~\ref{lem:orders} supplies a growth order; hence
\eqref{eq:minor-hypothesis} at $l=0$ gives $\det\Lmat[B]\neq0$.
Fix a child $w$ of $r$ and use the notation $B'=\{r,w\}$, $f'$ and
$M'=M_{f'}$ of Lemmas~\ref{lem:orders} and~\ref{lem:forest-fold}.
The bijection of Lemma~\ref{lem:orders} and \eqref{eq:weight-peel},
with the orders $\tau'$ grouped by $x=f(\tau'_1)\in B'$, give
\begin{equation}\label{eq:single-root-peel}
\sum_{\substack{\tau\in\Ord_B(f)\\ \tau_1=w}}\Pi_{B,f}(\tau)
=\sum_{x\in B'}\gam^{B,w}_x\,\mathcal G_x(B',f').
\end{equation}
By Lemma~\ref{lem:orders}, $f'$ is a $B'$-parent map, $B'$ is
proper because $k\geq2$, the minor condition
\eqref{eq:minor-hypothesis} holds for $(B',f')$, and
$\det\Lmat[B']\neq0$. Theorem~\ref{thm:order-mass} for the
two-element root set $B'$ gives
$\mathcal G_x(B',f')=\cof_x\bigl(\Nmat(B',M')\bigr)$ for
$x\in B'$; these cofactors are determinants of $1\times1$
matrices. That instance is proved by the induction on root sets
with at least two elements, so the argument is not circular.

The matrix $M'$ has unit row sums by Lemma~\ref{lem:blocks}, and
$|B|=1\leq n-2$. Lemma~\ref{lem:single-base} with $U=B$ and
$U'=B'$ evaluates the right side of \eqref{eq:single-root-peel}
as $\gam^B(\widetilde s+\mathbf e_w)_r$, and
Lemma~\ref{lem:forest-fold} gives
$\widetilde s+\mathbf e_w=\mathbf1_{D_w}$. Hence
\[
\sum_{\substack{\tau\in\Ord_B(f)\\ \tau_1=w}}\Pi_{B,f}(\tau)
=\gam^B(\mathbf1_{D_w})_r .
\]
Sum over the children $w$ of $r$. By Lemma~\ref{lem:blocks} the
sets $D_w$ partition $Z_r$, and $Z_r=Z$ because $r$ is the only
root. Linearity and \eqref{eq:unity} in Lemma~\ref{lem:harmonic}
give
\[
\sum_{\tau\in\Ord_B(f)}\Pi_{B,f}(\tau)
=\gam^B(\mathbf1_Z)_r=(\one_V)_r=1 .
\]
Under the Leibniz value $\det(\varnothing)=1$ this equality is
Theorem~\ref{thm:order-mass} at $B=\{r\}$; the proof above does
not rely on that reading.
\end{proof}

\section{Finite chains and confined excursions}

\begin{definition}[Chain, reversal, and first entrances]
\label{def:chain}
A stochastic matrix $P\in\R^{V\times V}$ has $P_{uv}\geq0$ and
$P\one_V=\one_V$. It is irreducible if, for every $u,v\in V$,
some $t\geq0$ has $(P^t)_{uv}>0$.
A stationary distribution is a vector $\pi\in\R^V$ with
$\pi_v>0$, $\sum_v\pi_v=1$, and $\pi^{\mathsf T}P=\pi^{\mathsf T}$.
For every stochastic $P$ put
$\lambda_v=\det\bigl((\Id_V-P)[V\setminus\{v\}]\bigr)$ and
$\Lambda=\sum_{v\in V}\lambda_v$. Given a stationary $\pi$, put
\begin{equation}\label{eq:reverse}
D=\operatorname{diag}(\pi),\qquad
P^*=D^{-1}P^{\mathsf T}D,\qquad
\Lmat=\Id_V-P^*.
\end{equation}
Let $(X_t)_{t\geq0}$ be the chain with transition matrix $P$
started at $r$. Set
$T_v=\inf\{t\geq0:X_t=v\}$ and
$T_{\mathrm{cov}}=\max_{v\in V}T_v$.
On the event of finite cover, its first-entrance tree has edges
$(X_{T_v-1},v)$ for $v\neq r$.
For $B\subsetneq V$ containing $r$, put $Z=V\setminus B$ and
$T_B=\max_{v\in Z}T_v$. On the same event, its stopped
first-entrance forest $F_B$ has edges $(X_{T_v-1},v)$ for $v\in Z$.
For a $B$-parent map $f$, write
\begin{equation}\label{eq:forest-weight}
w_P(f):=\prod_{v\in Z}\pi_{f(v)}P_{f(v)v}.
\end{equation}
\end{definition}

\begin{lemma}[Proper confined minors and Green series]
\label{lem:confined}
For every irreducible stochastic $P$ and every nonempty proper
$U\subset V$,
\begin{equation}\label{eq:confined}
\det(\Id_U-P[U])>0,
\qquad
G_U:=(\Id_U-P[U])^{-1}=\sum_{t\geq0}P[U]^t.
\end{equation}
There is $0\leq c_U<1$ such that, for every $m\in\mathbb N$,
every row sum of $P[U]^{m(n-1)}$ is at most $c_U^m$.
\end{lemma}
\begin{proof}
For each $i\in U$, irreducibility supplies a simple
positive-probability path from $i$ to $V\setminus U$ of length
at most $n-1$. Hence every row sum of $P[U]^{n-1}$ is less
than one. Finiteness of $U$ gives a common bound
$0\leq c_U<1$, since row sums are nonnegative;
submultiplicativity of the maximum row-sum norm proves the
geometric estimate. The Neumann series converges and multiplies
$\Id_U-P[U]$ to $\Id_U$.
The same bound applies to $tP[U]$ for $0\leq t\leq1$, so
$\det(\Id_U-tP[U])$ never vanishes. It is continuous and equals
$1$ at $t=0$; therefore it is positive at $t=1$.
\end{proof}

\begin{lemma}[Stationary distribution from principal cofactors]
\label{lem:stationary}
Every finite irreducible stochastic $P$ has a unique stationary
distribution. With $\lambda$ and $\Lambda$ as in
Definition~\ref{def:chain}, it is $\pi=\lambda/\Lambda$; in particular
$\lambda_v>0$ and $\Lambda>0$.
\end{lemma}
\begin{proof}
The positivity of $\lambda_v$ follows from
Lemma~\ref{lem:confined} for $U=V\setminus\{v\}$;
thus $\Lambda>0$.
Apply Lemma~\ref{lem:cofactor} to $\Id_V-P$, whose row sums
vanish. It gives
$\lambda^{\mathsf T}(\Id_V-P)=0$, so
$\pi=\lambda/\Lambda$ is a positive normalized stationary
distribution. A nonzero $(n-1)$-minor makes
$\operatorname{rank}(\Id_V-P)=n-1$. Its left kernel is therefore
one-dimensional, so normalization makes $\pi$ unique.
\end{proof}

\begin{lemma}[Reversal and normalized inverse columns]
\label{lem:green}
Let $P$ be irreducible and stochastic, and let $\pi$ be stationary.
Then $P^*$ is stochastic, $\Lmat\one_V=0$, and for every nonempty
proper $U\subset V$,
\begin{equation}\label{eq:green}
\det\Lmat[U]=\det(\Id_U-P[U])>0
\end{equation}
and
\begin{equation}\label{eq:green-reversal}
(G_U)_{ij}
=\frac{\pi_j}{\pi_i}(\Lmat[U]^{-1})_{ji}
\quad(i,j\in U).
\end{equation}
If $|U|\geq2$ and $j\in U$,
\begin{equation}\label{eq:inverse-ratio}
(\Lmat[U]^{-1})_{jj}
=\frac{\det\Lmat[U\setminus\{j\}]}{\det\Lmat[U]},
\qquad
\frac{\Lmat[U]^{-1}\mathbf e_j}
     {(\Lmat[U]^{-1})_{jj}}
=\gam^{U\setminus\{j\},j}|_U.
\end{equation}
\end{lemma}
\begin{proof}
Stationarity gives
$(P^*\one_V)_i=\pi_i^{-1}\sum_jP_{ji}\pi_j=1$.
Every entry $P^*_{ij}=\pi_jP_{ji}/\pi_i$ is nonnegative.
Also
$\Lmat[U]=D[U]^{-1}(\Id_U-P[U])^{\mathsf T}D[U]$;
its determinant and inverse give the equal determinant and
\eqref{eq:green-reversal}; the positivity in
\eqref{eq:green} follows from Lemma~\ref{lem:confined}.

The first identity in \eqref{eq:inverse-ratio} is Cramer's rule;
its numerator is positive by Lemma~\ref{lem:confined}, since
$U\setminus\{j\}$ is nonempty and proper.
Extend the inverse column $\Lmat[U]^{-1}\mathbf e_j$ by zero
off $U$ and divide by its positive diagonal entry. It is one
at $j$, zero outside $U$, and annihilated by the rows
$U\setminus\{j\}$ of $\Lmat$. Lemma~\ref{lem:harmonic}
identifies it with the stated harmonic extension.
\end{proof}

\begin{lemma}[Finite cover and first-entrance forest]
\label{lem:cover}
For an irreducible finite chain, every $T_v$ and
$T_{\mathrm{cov}}$ is finite almost surely. For a nonempty proper
$B\subset V$ with $r\in B$, the
first-entrance edges of $F_B$ form a $B$-parent map, and the
first-entrance tree is an arborescence rooted at $r$.
\end{lemma}
\begin{proof}
Apply the uniform escape estimate of
Lemma~\ref{lem:confined} to $U=V\setminus\{v\}$.
The chance of avoiding $v$ for $m(n-1)$ further steps decays
geometrically, from any current state outside $v$. Thus
$T_v<\infty$ almost surely; finiteness of $V$ gives the cover
claim. If $v\notin B$ is entered from $u$, then $T_u<T_v$.
Following its parent edges therefore strictly decreases first
entrance times, so it cannot cycle and must end at a vertex of
$B$. With $B=\{r\}$ this is an arborescence.
\end{proof}

\begin{lemma}[Confined-excursion formula]
\label{lem:excursion}
Under Definition~\ref{def:chain}, let $P$ be irreducible, let
$B\subsetneq V$ contain $r$, and let $f$ be a $B$-parent map.
For a growth order $\tau=(\tau_1,\ldots,\tau_k)$ put
$\tau_0=r$ and $U_l=U_l^\tau$. Then
\begin{equation}\label{eq:excursion}
\mathbb P_r(F_B=f,\ \text{discovery order}=\tau)
=\left(\prod_{v\in Z}P_{f(v)v}\right)
 \prod_{l=1}^{k}(G_{U_{l-1}})_{\tau_{l-1},f(\tau_l)}.
\end{equation}
For an ordering that is not a growth order, the probability is zero.
\end{lemma}
\begin{proof}
Immediately before the first entry into $\tau_l$, the visited
non-roots are $\tau_1,\ldots,\tau_{l-1}$, while the walk may
visit any root in $B$. It starts this interval at $\tau_{l-1}$
(at $r$ if $l=1$), spends $t_l\geq0$ steps within $U_{l-1}$,
then takes the edge $f(\tau_l)\to\tau_l$.
For fixed $t_l$ this interval has weight
$(P[U_{l-1}]^{t_l})_{\tau_{l-1},f(\tau_l)}
P_{f(\tau_l),\tau_l}$.
On finite cover, the successive intervals and their lengths are
uniquely determined; conversely, their concatenation gives exactly
the stated event. The cases indexed by the length vector are
disjoint. The Markov property makes their weights multiply, and
nonnegativity permits summing over all length vectors as a product
of the separate sums.
Summing each length with \eqref{eq:confined} gives
\eqref{eq:excursion}. Every recorded parent has already been
visited, so an incompatible ordering has probability zero.
\end{proof}

\begin{lemma}[Reversal and determinant telescoping]
\label{lem:joint-order}
Under Lemma~\ref{lem:excursion}, let $\pi$ be the stationary
distribution of $P$ and use the forest weight
\eqref{eq:forest-weight}.
Then for each $\tau\in\Ord_B(f)$,
\begin{equation}\label{eq:joint-order}
\mathbb P_r(F_B=f,\ \text{discovery order}=\tau)
=\frac{\det\Lmat[B]}
 {\pi_r(\prod_{v\in Z}\pi_v)\Lambda}\,
 w_P(f)\,(\Lmat[B]^{-1})_{f(\tau_1),r}
 \Pi_{B,f}(\tau).
\end{equation}
The formula includes $k=1$, when $\Pi_{B,f}(\tau)=1$.
\end{lemma}
\begin{proof}
Apply \eqref{eq:green-reversal} to each Green entry of
\eqref{eq:excursion}. Since
$\prod_{l=1}^k\pi_{\tau_{l-1}}
=\pi_r(\prod_{v\in Z}\pi_v)/\pi_{\tau_k}$, the factors
outside the inverse entries become
$w_P(f)\pi_{\tau_k}/[\pi_r\prod_{v\in Z}\pi_v]$.

For $2\leq l\leq k$ apply \eqref{eq:inverse-ratio} to
$U_{l-1}=U_{l-2}\cup\{\tau_{l-1}\}$ and
$j=\tau_{l-1}$. It gives
\[
(\Lmat[U_{l-1}]^{-1})_{f(\tau_l),\tau_{l-1}}
=\gam^{U_{l-2},\tau_{l-1}}_{f(\tau_l)}
  \frac{\det\Lmat[U_{l-2}]}{\det\Lmat[U_{l-1}]}.
\]
The first inverse entry remains
$(\Lmat[B]^{-1})_{f(\tau_1),r}$.
The other determinant ratios telescope to
$\det\Lmat[B]/\det\Lmat[U_{k-1}]$; this quotient is $1$
when $k=1$. Since
$U_{k-1}=V\setminus\{\tau_k\}$,
\eqref{eq:green} and Lemma~\ref{lem:stationary} give
$\det\Lmat[U_{k-1}]=\lambda_{\tau_k}
=\pi_{\tau_k}\Lambda$.
The remaining harmonic factors form $\Pi_{B,f}(\tau)$, proving
\eqref{eq:joint-order}.
\end{proof}

\section{Aldous--Broder}

\begin{definition}[Tree and forest families]
\label{def:families}
For nonempty $B\subseteq V$, let $\mathcal F_B$ be the finite
set of $B$-parent maps, identified with their edge sets
$\{(f(v),v):v\in V\setminus B\}$. If $B=\{r\}$,
these are the arborescences rooted at $r$.
\end{definition}

\begin{theorem}[Aldous--Broder for an irreducible chain]
\label{thm:ab}
Let $P$ be irreducible and stochastic, let $\pi$ be its stationary
distribution, and start the covering walk at $r\in V$.
For every $f\in\mathcal F_{\{r\}}$, its first-entrance tree $F$
satisfies
\begin{equation}\label{eq:ab}
\mathbb P_r(F=f)
=\frac{\displaystyle\prod_{v\neq r}\pi_{f(v)}P_{f(v)v}}
 {\displaystyle\Lambda\prod_{v\in V}\pi_v}
=\frac{\displaystyle\prod_{v\neq r}\pi_{f(v)}P_{f(v)v}}
 {\displaystyle\sum_{f'\in\mathcal F_{\{r\}}}
   \prod_{v\neq r}\pi_{f'(v)}P_{f'(v)v}}.
\end{equation}
\end{theorem}
\begin{proof}
In Lemma~\ref{lem:joint-order} take $B=\{r\}$.
Every growth order starts with a vertex whose parent is $r$,
and the one-by-one inverse entry is
$(\Lmat[\{r\}]^{-1})_{rr}=1/\det\Lmat[\{r\}]$.
Thus, for each compatible discovery order $\tau$,
\[
\mathbb P_r(F=f,\ \text{order}=\tau)
=\frac{w_P(f)}{\Lambda\prod_{v\in V}\pi_v}
  \Pi_{\{r\},f}(\tau).
\]
Summing over growth orders and applying
Theorem~\ref{thm:single-root} gives the first expression in
\eqref{eq:ab}. By Lemma~\ref{lem:cover}, the first-entrance
tree belongs to exactly one member of $\mathcal F_{\{r\}}$
almost surely. Sum the first expression over this finite family;
the resulting equality $1=\sum_{f'}w_P(f')/
(\Lambda\prod_v\pi_v)$ identifies its denominator with the
tree-weight sum, giving the second expression.
\end{proof}

\section{Stopped multi-root forest law}

\begin{definition}[Order masses and completion factor]
\label{def:completion}
Let $P$ be irreducible and stochastic with stationary
distribution $\pi$, let $B\subsetneq V$ be nonempty, $r\in B$, and
$f\in\mathcal F_B$. With the matrix $\Lmat$ in
\eqref{eq:reverse}, write
\[
g(B,f):=(\mathcal G_i(B,f))_{i\in B},\qquad
\varphi_r(f):=g(B,f)^{\mathsf T}\Lmat[B]^{-1}\mathbf e_r.
\]
\end{definition}

\begin{lemma}[Forest as a marginal of the covering tree]
\label{lem:forest-marginal}
Let $P$ be irreducible and stochastic with stationary
distribution $\pi$, let $B\subsetneq V$ be nonempty, choose
$r\in B$, and let $f\in\mathcal F_B$. Define
\[
\mathcal C_r(f)
:=\{h\in\mathcal F_{\{r\}}:
   h|_{V\setminus B}=f\}.
\]
Then
\begin{equation}\label{eq:forest-marginal}
\mathbb P_r(F_B=f)
=\sum_{h\in\mathcal C_r(f)}\mathbb P_r(F=h)
=\frac{\displaystyle\sum_{h\in\mathcal C_r(f)}w_P(h)}
 {\displaystyle\Lambda\prod_{v\in V}\pi_v}.
\end{equation}
For $h\in\mathcal C_r(f)$,
$w_P(h)=w_P(f)\prod_{j\in B\setminus\{r\}}
\pi_{h(j)}P_{h(j)j}$.
\end{lemma}
\begin{proof}
The first-entrance edge into any $v\in Z$ is fixed at time
$T_v\leq T_B\leq T_{\mathrm{cov}}$. Thus the stopped forest is
exactly the restriction of the full first-entrance tree to edges
entering $Z$. The disjoint events $\{F=h\}$ with
$h\in\mathcal C_r(f)$ partition $\{F_B=f\}$.
Apply Theorem~\ref{thm:ab} to each event. The factorization of
$w_P(h)$ separates its edges entering $Z$ from those entering
$B\setminus\{r\}$.
\end{proof}

\begin{theorem}[Stopped forest law and completion weights]
\label{thm:stopped}
Let $P$ be irreducible and stochastic with stationary
distribution $\pi$, let $B\subsetneq V$ be nonempty, choose a
seed $r\in B$, and let $f\in\mathcal F_B$. Then
\begin{equation}\label{eq:stopped}
\mathbb P_r(F_B=f)
=\frac{\det\Lmat[B]}
 {\pi_r(\prod_{v\in Z}\pi_v)\Lambda}
\,w_P(f)\,\varphi_r(f).
\end{equation}
The total weight of the arborescence completions of $f$ is
\begin{equation}\label{eq:completion-weight}
\sum_{h\in\mathcal C_r(f)}
 \prod_{j\in B\setminus\{r\}}
 \pi_{h(j)}P_{h(j)j}
=\det\Lmat[B]
 \left(\prod_{j\in B\setminus\{r\}}\pi_j\right)
 \varphi_r(f).
\end{equation}
\end{theorem}
\begin{proof}
Sum \eqref{eq:joint-order} over the compatible growth orders,
grouping them by the root $i=f(\tau_1)$. The common factor is
$\det\Lmat[B]w_P(f)/
[\pi_r(\prod_{v\in Z}\pi_v)\Lambda]$;
the remaining sum is
\[
\sum_{i\in B}\mathcal G_i(B,f)
   (\Lmat[B]^{-1})_{ir}
=g(B,f)^{\mathsf T}\Lmat[B]^{-1}\mathbf e_r
=\varphi_r(f).
\]
This proves \eqref{eq:stopped}, including when $w_P(f)=0$.

If $w_P(f)>0$, compare \eqref{eq:stopped} with
\eqref{eq:forest-marginal}, use the factorization of $w_P(h)$,
and cancel $w_P(f)$. Since
$\prod_{v\in V}\pi_v=
(\prod_{v\in Z}\pi_v)(\prod_{j\in B}\pi_j)$,
the resulting equality is \eqref{eq:completion-weight}.

For arbitrary $f$, set
$P_\varepsilon=(1-\varepsilon)P+
\varepsilon\one_V\one_V^{\mathsf T}/n$, with
$0<\varepsilon<1$. All its entries and hence all its forest
weights are positive, so \eqref{eq:completion-weight} holds for
$P_\varepsilon$. By Lemma~\ref{lem:stationary}, its stationary
distribution is the normalized cofactor vector of
$\Id_V-P_\varepsilon$, and therefore converges to $\pi$ as
$\varepsilon\downarrow0$. Lemma~\ref{lem:green} ensures that
all principal minors used in the harmonic extensions and in
$\Lmat[B]^{-1}$ are positive at the limit; continuity keeps them
nonzero for sufficiently small $\varepsilon$. The set of growth
orders is finite, so its order masses and $\varphi_r(f)$ are
continuous at the limit. Every term on both sides of
\eqref{eq:completion-weight} is thus continuous there.
Taking the limit proves the formula also when $w_P(f)=0$.
\end{proof}

\section{Mathematical contract for the Lean development}

The three public endpoints are Theorems~\ref{thm:order-mass},
\ref{thm:ab}, and~\ref{thm:stopped}, including both displayed
conclusions of Theorem~\ref{thm:stopped}. Theorem~\ref{thm:single-root}
and the lemmas before each endpoint are its proof obligations, not
additional application targets. In Sections 1--4, $\Lmat$ is an
arbitrary real zero-row-sum matrix where stated. In the chain
sections it is specialized to $\Id-P^*$ by \eqref{eq:reverse}.

The following boundary conditions are part of the statements to be
formalized.
\begin{enumerate}
\item A root set in the order--mass theorem is nonempty and proper,
so $k=|V\setminus B|\geq1$. The peeling bijection and the fold are
used only for $k\geq2$; the $k=1$ case is proved directly. No
``growth order'' for $B=V$ is silently used.
\item Theorem~\ref{thm:order-mass} needs nonzero minors only for
proper bases $U_l^\tau$ of growth orders of the \emph{given}
parent map. Neither all intermediate subsets nor positivity of
the matrix entries is assumed. The conclusion holds for every
root, including one with no children. Only its reading at $|B|=1$
takes a cofactor of a $1\times1$ matrix, with the Leibniz value
$\det(\varnothing)=1$. Theorem~\ref{thm:single-root} is proved
through Lemma~\ref{lem:single-base} and partition of unity, not
by that reading; the Aldous--Broder and stopped-forest endpoints
depend on it only through Theorem~\ref{thm:single-root}.
\item The Aldous--Broder and stopped-forest statements assume an
irreducible row-stochastic matrix, not reversibility,
aperiodicity, or strictly positive transitions. Their probabilities
refer to the Markov chain started at the stated seed. The
first-entrance edge of a non-root $v$ is $(X_{T_v-1},v)$; the
seed $r$ is never assigned such an edge.
\item Equation~\eqref{eq:completion-weight} is asserted even when
$w_P(f)=0$. It cannot be obtained by cancelling that weight.
Its separate perturbation-and-continuity argument is required.
\end{enumerate}

The manuscript's $B$-parent map is a function on $Z=V\setminus B$.
A total Lean function $f:V\to V$ may be used internally only after
proving that values on $B$ do not affect the forest, orders, blocks,
weights, or theorem statement. The paper-facing theorem should
accept a map on the non-roots, or provide an explicit equivalence
with that formulation. Likewise, $\cof_i(A)$ always deletes both
row and column $i$; the two matrix levels $\Lmat[B]$ and
$\Nmat(B,M_f)$ must not be conflated.


\end{document}
